\chapter{Social}

\begin{multicols}{2}


\section{Social Actions:}

Social actions are skill tests between persuasion, deception, will, insight, cunning and knowledge. Which skill is used is determined by the context of the action and the GM's decision. 

\textbf{How to do it?:} Social actions are performed through character interpretation. A character says or does something that leads a social consequence. If the social action is performed within a combat context, the GM may decide to require a certain amount of Action Points (AP), usually 3, to perform it. It could be less for short sentences and more for speeches or explanations.

\textbf{Defender rolls:} Social actions are always resisted by the defender. The attacker tells them their skill value, the IB gets added to that, and the defender adds their own skill values and the dice result. The degree of success does not need to be disclosed to the other players. This means that the attacker will have to interpret the defender's reaction to determine how well they did. The defender can also choose to give a false impression of their degree of success, like pretending to be afraid, or to believe in something without the attacker knowing. This will require the defender to be honest about their degree of success, so this mechanic may not work well with some players. In that case, just switch to a more transparent system where the degree of success is disclosed.

\textbf{Interpretation Bonus (IB):} Interpretation is a very important factor in determining the effectiveness of social actions. If the player's interpretation in the scene is not suitable for the social action they are trying to perform, a negative IB should be applied. The action may even be denied by the GM. In cases where the probability of the social action succeeding is very high and the interpretation has been very good, a high positive IB may be granted. The bonuses and penalties should range between -10 to +10. Larger values can be used instead of a straight veto or auto pass.

\textbf{Choosing an IB:} 

\begin{itemize}
  \item An IB of 10 is almost a guaranteed success, it should be given only in situations where the NPC should want to agree with the player anyways and the acting was really good and cohesive with the scene. Often, it is better to just allow the player to skip the test and get the desired outcome, but if the GM wants to keep the test, this is the IB to use.
  \item IBs around +5 are the most common ones for situations where a character is very likely to agree unless there is some special reason not to. A failure in such a test is just a prompt to find such a special condition to failure. Use this when player has made a very good argument and good acting. They are a reward for good roleplay, reading the scene well and coming up with practical game solutions.
  \item IBs between +2 and -2 are for common situations, usually improvised, where the one character does not have a strong opinion on the matter. An example would be to ask for a free pint of beer in a tavern. Give a bonus or penalty to based on how invested the player is in the interaction and how they show that in roleplay.
  \item IBs around -5 are for situations where the player is trying to do something that the NPC does not want to do, but it is not completely out of the question. For example, asking a shopkeeper for a discount out of the blue may warrant an IB of -5, but giving them a descent reason to do it may reduce the penalty to -3 or lower. This is a penalty to someone forcing the action where it does not belong or when the player has said something really bad in the roleplay, like insulting the NPC or being completely out of character. A success in such an action is a prompt to come up with a good reason why that should work.
  \item Finally, -10 is just to avoid giving straight vetos. It's an answer to at least let them try to do it.
\end{itemize}
  

\textbf{Test Outcomes:} The outcome of a social action depends on what skill is used for the test and the context of the action. An intimidation that goes wrong will have worse consequences than an attempt to seduce someone by giving a gift. The consequences should be aligned with the degree of success of failure. Maybe a critical success includes a bonus, or maybe a critical failure makes the NPC hostile. The consequences should be determined by the GM, but they should be consistent with the context and the degree of success or failure.

\textbf{Negotiating, Convincing, and Explaining:} These are the most common social actions. They are all persuasion tests, usually against insight or knowledge, depending on the what has been said and what kind of cue the defender is looking for. If they want to focus on the vibes to get a sense of the other person's intentions, the will use insight. If they want to check the logic of the argument, they will use knowledge. The description of what is descovered from a successful defense should be different for each defense.  

\textbf{Lying and falsehoods:} When anyone gives an information, an insight vs deception check can be called to determine whether the information is believed or not. Anyone can call for that test, but when the GM does, it must be done. On a crit or hit, the one who used insight knows the other one's true intentions, but not necessarily the truth of what they said. On a miss or graze, the one who used insight is convinced that the other one is telling the truth, even if they are lying. 

\textbf{Persuasion vs Deception:} Sometimes it can be ambiguous whether a social action is persuasion or deception. As a rule of thumb, the more important the lie is for the success of an interaction, the greater the indication that deception should be used. On the other hand, if the focus is in the overall story, feelings, or the relationship between the characters, persuasion is more appropriate. For example, if you are trying to convince someone to give you money by telling them that you need it to pay for your sick mother's medicine, but none of that is true, that is a deception. If it were true, it would be persuasion. But ybe they don't care about your mother, and they will just say that they are sorry and move on, even in a critical success. In that case, you have to pass two tests, but use only one dice roll for both. The first test is to see if you can convince them that you are telling the truth, and the second test is to see if they care about what you are saying.  

\textbf{Changing emotional state:} Social actions can be used to make a character afraid, enraged, or distracted. The GM should determine the appropriate skill to use for each case. 

\textbf{Intimidation:} The usual intimidation attempt is meant to make someone afraid and is made as a will vs will test. For example, when you look someone in the eye with murderous intent, you can use will vs will to make them afraid. Explaining all the ways someone's business is going to fail and telling them just how they are going to be destitute is an intimidation that uses persuasion against will, insight or even knowledge, depending on whether the defender wants to rely on their will to fight, their intuition that the other is bluffing or their knowledge of the situation. The IB should reflect the physical reality of the situation. A giant armed muscled and armored barbarian is obviously scarier than a 10 year old kid. But a dragon probably finds both equally  This difference should be observed in the IB.

\textbf{Provocation:} When you insult someone, you can use persuasion vs insight or will to make them enraged. As a recommendation, insults are most effective when they exploit someone's frustrations or pride. Subtelty and timing are important. For example, if you know that someone is very proud of their appearance, insulting their looks may be a good way to provoke them. That is especially effective if you tell them they got rejected romantically because fo their big nose. 

\textbf{Seduction:} When you seduce someone, you can use persuasion vs insight or will to make them charmed. The best way to seduce someone is to make them feel good about themselves and to make them feel a connection with you. Compliments, romantic gestures, and showing vulnerability can be effective ways to seduce someone. For example, if you know that someone is very insecure about their intelligence, complimenting their intellect and telling them how much you admire it can be a good way to seduce them. The more subtle and the better the timing, the better the IB.

\textbf{Distraction:} When distracting someone, you can use persuasion vs cunning in a situation where you are just making a little show, or telling someone a story to divert attention. 

\textbf{When to allow Tests and Repetition:} it is important to make every social test count to the end outcome of the interaction. This means that there must be irreversible consequences to the outcome of a test. The idea behind this is to avoid allowing multiple tests to evaluate the same intent. For example, if you try to intimidate someone by staring them down, you cannot attempt to do that later again, nor can anyone else in your group, unless their circumstances are deemed very different. In the same note, if that test succeeds, the target is now afraid for the rest of the interaction.

\textbf{Group tests:} are made with the skill of a single character. The other ones involved add to the IB, depending on how they change the scene, but their social skills are not used in the test.

\textbf{Degrees of success in social interactions:} Think about a hit in a social test as a success in the social action's intent. The attacker gets what they wanted, no more. A critical is meant to be a success and more. A critical intimidation may be too effective and make someone run away or have a panic attack instead of telling what you want to know. They may tell you more than what you asked for. A graze is a basic failure. It is a standard no. You don't get what you wanted and you burn a chance to get it, your objective should take a setback but the defender will not overreact and shut down the interaction. A miss is a failure and more. The defender will see through your intent, lose respect in you, lose trust, take precautions against your intent and maybe even turn hostile. A miss can ruin the entire interaction.

\textbf{Using skills to skip tedious social procedures:} Sometimes players may not want to go through the roleplay of a social scene. In these cases, they can use usually persuasion but sometimes other appropriate skill to represent their social skill and skip the scene and go to the outcome. The player should give a general description of what they are trying to achieve, and the GM should determine the IB based on how well it makes sense. A hit or crit means that the desired outcome is achieved.

\section{Contacts}

Contacts provide information, utilities, and services. Those who are not contacts usually charge for exclusive information, charge higher prices for the same goods, hide special stock.\\

Any social action that requires someone to be a contact gets an immediate -5 penalty if they are not one. This means that the penalty can only be removed through roleplay, by convincing the character to treat you like a contact. \\

The types of contacts and what they do are described below: \\

\textbf{Artisan:}
Any artisan who is not a contact charges +50\% on the value of any item by default. \\
Blacksmith: sells and repairs weapons and armor. may accept later payments as a contact. \\
Alchemist: sells potions and buys ingredients. 
Tailor: Makes and adjusts clothes. 

\textbf{Merchant}
Any merchant who is not a contact charges +50\% on the value of any item by default. \\
-merchant: Buys and sells common goods.
-smuggler: Gets illegal items. Smuggles things into places.
-fence: buys rare objects for a fair price.
-tamer: sells horses, dogs, and other trained animals.

\textbf{Service}
Worker: performs construction work
Mercenary
clergyman: does various tasks
doctor: heals injuries

\textbf{Influencer}:
Bard - increases popularity
Noble: grants access to places and can impress vassals
Loan: lends money, offers work
spy: knows information about a certain place

\textbf{Trainer}:
Martial
Knowledge
Spirit
Constitution

\textbf{Acquiring contacts:} This can be done by completing missions, or by making an acting test to see how much time and money will be spent on acquiring this contact. The amount of time and money required depends on the relevance and utility of the contact.

\section{Tags}

Tags are given to NPCs to represent their desires, biases, interests and overall social inclinations. They can be used by other characters as a means to approach this NPC and get them to treat you better temporarily. Some tags can be very strong and make an NPC treat you like a contact, while others can be just a bonus to subsequent social actions.

Some tags may require a skill test to be exploited. If an NPC is affiliation bound, but you don't belong, you may have to make a deception check against their insight to be able to pass for a member of the same clan.

An insight test can be made to try to identify an NPCs tags. Like othe social actions, insight is liable to an IB, where if the NPC has nothing to demonstrate, the test becomes harder. Once the player uses insight once, they can only do it again against the same character if circumstances change significantly and the GM calls for it. The player is only allowed to try insight by their initiative once.

Tags are can be temporary features or NPC characteristics. A town guard may act status oriented when on duty. An attempt to enter a guarded room will succeed automatically if you have the right credentials, no matter the NPCs tags.

Tags are grouped into a few big categories to make it easier to improvise. The GM can decide quickly to use one of the big categories on an NPC and come up with a specific tag in the middle of the interaction.

\textbf{1. Status-Oriented}
Responds to:

Reputation (Respect): Being recognized for your good reputation. High popularity may waive the need for the test. \\
Wealth: Attempts  to convince someone of your wealth. The test may be waived if there is sufficient physical proof. \\
Titles \\
Displays of importance \\
“Is this person above or below me?”

\textbf{2. Affiliation-Bound}
Responds to:

Shared contacts: Having a relevant common contact with someone may be enough to gain their trust. \\
Factions: \\
Group loyalty \\
Mutation: Many characters reject mutants. \\
“Are you one of us?”

\textbf{3. Usefulness-Oriented}
Responds to:

Competence: \\
Skill displays: Some characters need competent and useful people. The impress test is necessary when their skills are not yet recognized. \\
Practical solutions \\
Proven results \\
“Can you solve my problem?”

\textbf{4. Seduceable}
Some characters trust those who seduce them.
Responds to:

Beauty \\
Charm \\
Romantic intent \\
Vulnerability cues \\
\textbf{key question:} “Do I want you?” \\
These NPCs are looking for human connection or intimacy. 

\textbf{5. Order-Bound}
Responds to:

Etiquette: Some environments require appropriate behavior. It is necessary to have knowledge of the etiquette of a certain place to be treated well. \\
Ritual \\
Proper conduct \\
Cultural fluency \\
“Are you behaving correctly?”

\textbf{6. Threat-Sensitive}
Responds to:

Stress \\
Violence cues \\
Power imbalance \\
Consequences \\
Notoriety: Reputation for being a liar, malicious, or incompetent has a negative impact on any negotiation. \\
“What happens if I refuse?”

\textbf{7. Principle Oriented}
Responds to:

Ideology \\
Morality \\
Honor \\
“What do you believe in?”

\textbf{8. Resource Oriented}
Responds to:

Wealth \\
Access \\
Service \\
“What can you give me?”

\section{Reputation}

There are 3 forms of reputation: respect, fear, and popularity. Each NPC has a different level of reputation required to affect them.

\textbf{Respect:} +3 Charm. ranges from 1-10
Respect means that the character is considered more competent, trustworthy, and important. Each NPC has a level of respect requirement. If the tolerance is equal to or lower, the respect bonuses apply. \\
Characters with high respect will be better paid, will be called for more important tasks, and will have access to more exclusive areas of society.

\textbf{Notoriety:} +3 Stress. ranges from 1-10
Notoriety is the reputation for being cruel, aggressive, or a liar. Each character has a tolerance for notoriety. Bonuses apply if the tolerance is lower than the notoriety. \\
Characters with high notoriety will be called for immoral or illegal tasks.

\textbf{Popularity:}  Ranges from 1-10
Popularity is the measure of how many people know you. Characters with high popularity are recognized by everyone in the community and waive first impressions. Characters have a level of information. If the popularity is greater than the level of information, the bonuses apply.
Popular characters are called for diplomatic activities.

Reputation is local. It is necessary to have some contact in other places to transfer reputation.

Popularity declines monthly. It never naturally falls below 4.

\subsection{Increasing Reputation}

\textbf{Contacts:} Having important contacts contributes to reputation.

\textbf{Scandals, heroic deeds, completed missions:} Reputation is part of the reward for playing the game. The higher the reputation, the harder it is to increase.

\textbf{Celebrating:} It is possible to spend downtime improving popularity. The amount of money spent and the character's charisma determine how much reputation is gained.

\section{Factions}

\subsection{Characterization:}

a name:
a goal:
a belief:

key characters: Leader, manager, specialist, trainer 
Allies and enemies

\subsection{Military}
Power of the army. 
Quantity and quality of troops.

Maintaining troops:
Buying troops: 
Recruiting troops:

examples: Mercenary Organization, Royalty

\subsection{Wealth}
Ownership of strategic resources, high cash flow, large cash reserves.

Income:
Cost:
Reserve:

examples: Merchants, Mine Owner

\subsection{Influence:}
ability to change people's opinions.

Credibility: People believe what you say.
Popularity: People like you

examples: University, Church

\end{multicols}
